% ============================================================================
\clearpage
\section*{v4.9.0-r11 Source-Only: A Minimal Dynamics for Reading B --- Kicked Duality Rotation}
\addcontentsline{toc}{section}{v4.9.0-r11 Source-Only: A Minimal Dynamics for Reading B --- Kicked Duality Rotation}
\label{sec:v490-r11-kicked-duality-dynamics}
% ============================================================================

\paragraph{Purpose.}
r10 found that the open threads of the core meet at one gap: Reading~B has no dynamics (r10 of v4.8.0).  The C5 factor
\(3/4\) would be derived if a dynamics produced the alternation ``rotate, then compensate'' of r10.1.  A shape-(iii) witness
would need a dynamics that changes the invariant \(s\) (r10.2).  This layer proposes the smallest dynamics with these
properties and checks what it does.  It is a model proposal, not a derivation from the source.  The script is
\texttt{llm/k52\_kick\_dynamik.py}.  No value is compared with data.  Earlier layers are not rewritten.
This layer does not assert or approach an exact finite C5 zero closure.

\noindent\textbf{Source-only status.}
This v4.9.0-r11 layer is source-only.  It does not build a PDF, does not assign a new DOI,
and does not authorize a public release.  The v4.9.0 Version DOI
\href{https://doi.org/10.5281/zenodo.23534954}{10.5281/zenodo.23534954} and the Concept DOI
\href{https://doi.org/10.5281/zenodo.19210728}{10.5281/zenodo.19210728} are unchanged by it.

\subsection*{v4.9.0-r11 the proposal}
\label{subsec:v490-r11-proposal}

\paragraph{Definition (kicked duality dynamics, proposal).}
Between kicks the field turns in the duality circle at a constant rate, \(z(t)=e^{i\omega t}z\), on the doublet
\(F\mapsto FR_{\omega t}^{\mathsf T}\).  At the times \(t_n=n\tau\) the compensator \(C_{\HC}\) is applied, and the content of the
partner port is handed to a reservoir:
\[
F_{n+1}=F_nR_{\omega\tau}^{\mathsf T}C_{\HC},\qquad
\mathcal E_{\rm res}^{(n+1)}=\mathcal E_{\rm res}^{(n)}+\|F_nR_{\omega\tau}^{\mathsf T}A\|^2 .
\]
The rotation is the transport of Reading~B, the kick is the readout of the main body, and the reservoir is the partner
sector \(\operatorname{ran}A\) of r03.1.

\subsection*{v4.9.0-r11 what it does}
\label{subsec:v490-r11-properties}

\paragraph{Theoremlet v4.9.0-r11.1.}
\begin{enumerate}[leftmargin=2em]
\item[(a)] Each stage keeps \(\cos^2\omega\tau\) of the content and exports \(\sin^2\omega\tau\).  Retained plus exported
content is conserved.
\item[(b)] The stage factor equals \(3/4\) for exactly one \(\omega\tau\in(0,\pi/2)\), namely \(\omega\tau=\pi/6\).  Then the
retained content runs through \(3/4,9/16,27/64,81/256,\ldots\) and first falls below \(1/256\) at stage \(20\).
\item[(c)] In the state disk of r04 a stage is a rotation of the Bloch vector by \(2\omega\tau=60^\circ\) followed by a
measurement along the \(J\)-axis, kept with probability \((1+\cos60^\circ)/2=3/4\).
\item[(d)] The duality phase accumulates in steps of \(30^\circ\).  Three stages make the quarter turn from matter to
tachyon, six the half turn \(z\mapsto-z\) (the charge conjugation of r02.2), twelve the full turn, twenty-four the return of
the spinor of r06.4.
\item[(e)] The rotation leaves \(|s|\) unchanged.  The kick changes it: it sets \(\operatorname{Re}s=0\) and lowers
\(|s|\) (in the checked example from \(2.15\) to \(0.83\)).  The dynamics therefore leaves the transport group of
Reading~B, as r10.2 requires of any witness that changes \(s\).
\item[(f)] Dividing the same phase \(\pi/6\) into \(m\) kicks keeps \(\cos^{2m}(\pi/6m)\): \(0.75\), \(0.87\), \(0.95\), \(0.99\)
for \(m=1,2,5,50\).  The stage factor is set by how often the readout acts per phase, which is the quantum Zeno pattern.
\end{enumerate}

\paragraph{Proof sketch.}
(a)--(b) \(C_{\HC}R_\alpha C_{\HC}=\cos\alpha\,C_{\HC}\) (r10.1) and \(C_{\HC}+A=1\).  (c) \(G\mapsto RGR^{\mathsf T}\) turns the
Bloch vector by twice the angle (r04.2).  (e) direct computation.  All items were checked numerically, (b) also in exact
rational arithmetic.  \(\square\)

\subsection*{v4.9.0-r11 what it does not do}
\label{subsec:v490-r11-limits}

\begin{enumerate}[leftmargin=2em]
\item \emph{No finite closure.}  The retained content after \(n\) stages is \((3/4)^n>0\) for every \(n\).  It vanishes in
one stage only for \(\omega\tau=\pi/2\), and then the whole content sits in the reservoir: it is exported, not annihilated.
The dynamics reproduces \(C5_{\rm det/eff}\) at \(\theta=1/256\), \(n=20\), and keeps
\(0_{\rm termdata}\not\Rightarrow0_{\rm all}\).  The all-readout slot includes the reservoir.
\item \emph{No scale.}  Only the product \(\omega\tau\) is fixed.  The source supplies neither a rate \(\omega\) nor an interval
\(\tau\).  The oscillator frequency of Theorem~\ref{thm:spring}, \(\gstring\sqrt{2/3}\), is a dimensionless candidate for
\(\omega\) in the source's units.  Without a unit of time it fixes nothing.
\item \emph{No mechanism for the kick.}  Why the readout acts at regular intervals, and three times per quarter turn, is
not explained.  The triolic division of r10 is assumed, not derived.
\item \emph{Not the field equations.}  The proposal acts on the doublet at a point.  It is not a field theory and makes no
claim about the Maxwell or Einstein limits of the main body.
\end{enumerate}

\paragraph{Remark.}
The structure, a fixed polarisation rotation followed by a polariser at every cycle, is the one used in quantum Zeno
experiments with photons, among them the interaction-free measurements of Kwiat and co-workers (1995).  There a rotation of
\(\pi/(2N)\) per cycle with \(N\) cycles per quarter turn leaves \(\cos^{2N}(\pi/2N)\).  The proposal of this layer is that
structure with \(N=3\).  The correspondence concerns the algebra only.  On a light wave the compensator of r03 is not a
polariser (r03.2).

\subsection*{v4.9.0-r11 register entries}
\label{subsec:v490-r11-register}

\begin{center}
\small
\begin{tabular}{@{}>{\raggedright\arraybackslash}p{0.34\textwidth}>{\raggedright\arraybackslash}p{0.58\textwidth}@{}}
\toprule
statement & status after v4.9.0-r11\\
\midrule
dynamics of Reading~B (r10 of v4.8.0) & proposal: kicked duality rotation with compensator kicks and partner reservoir; model decision, not derived\\
C5 factor \(3/4\) & produced by the proposal at \(\omega\tau=\pi/6\), uniquely in \((0,\pi/2)\) (r11.1)\\
shape (iii) & the kick changes \(s\), but the dynamics exports instead of annihilating; no finite closure (r11)\\
scale of the dynamics & open: only \(\omega\tau\) fixed\\
\bottomrule
\end{tabular}
\end{center}

\subsection*{v4.9.0-r11 ledger and guard update}
\label{subsec:v490-r11-ledger-guard}

\[
R_{\rm dyn}^{v490\text{-}r11}=(0,0,1)
\]
for (a minimal dynamics proposed, its stage factor and cycles established, its scale and kick mechanism open).

The layer spends no guard:
\[
R_{\rm glob}^{(v490\text{-}r11,+)}=(0,1,1,1,1),
\qquad
\Gamma_{\rm guard}^{v490\text{-}r11}=(1,1,1,1,1),
\qquad
\beta_4=1/256\ \text{unspent}.
\]
It does not imply exact finite \(C5\) zero closure, exact all-readout terminality,
\(C6/T2\), a global photon-clock/RH theorem, or global all-readout closure.

